% Options for packages loaded elsewhere
\PassOptionsToPackage{unicode}{hyperref}
\PassOptionsToPackage{hyphens}{url}
\PassOptionsToPackage{dvipsnames,svgnames,x11names}{xcolor}
%
\documentclass[
  10pt,
  twocolumn,
]{article}
\usepackage{amsmath,amssymb}
\usepackage{setspace}
\usepackage{iftex}
\ifPDFTeX
  \usepackage[T1]{fontenc}
  \usepackage[utf8]{inputenc}
  \usepackage{textcomp} % provide euro and other symbols
\else % if luatex or xetex
  \usepackage{unicode-math} % this also loads fontspec
  \defaultfontfeatures{Scale=MatchLowercase}
  \defaultfontfeatures[\rmfamily]{Ligatures=TeX,Scale=1}
\fi
\usepackage{lmodern}
\ifPDFTeX\else
  % xetex/luatex font selection
\fi
% Use upquote if available, for straight quotes in verbatim environments
\IfFileExists{upquote.sty}{\usepackage{upquote}}{}
\IfFileExists{microtype.sty}{% use microtype if available
  \usepackage[]{microtype}
  \UseMicrotypeSet[protrusion]{basicmath} % disable protrusion for tt fonts
}{}
\makeatletter
\@ifundefined{KOMAClassName}{% if non-KOMA class
  \IfFileExists{parskip.sty}{%
    \usepackage{parskip}
  }{% else
    \setlength{\parindent}{0pt}
    \setlength{\parskip}{6pt plus 2pt minus 1pt}}
}{% if KOMA class
  \KOMAoptions{parskip=half}}
\makeatother
\usepackage{xcolor}
\usepackage[paper=a4paper,margin=0.8in]{geometry}
\usepackage{color}
\usepackage{fancyvrb}
\newcommand{\VerbBar}{|}
\newcommand{\VERB}{\Verb[commandchars=\\\{\}]}
\DefineVerbatimEnvironment{Highlighting}{Verbatim}{commandchars=\\\{\}}
% Add ',fontsize=\small' for more characters per line
\newenvironment{Shaded}{}{}
\newcommand{\AlertTok}[1]{\textcolor[rgb]{1.00,0.00,0.00}{\textbf{#1}}}
\newcommand{\AnnotationTok}[1]{\textcolor[rgb]{0.38,0.63,0.69}{\textbf{\textit{#1}}}}
\newcommand{\AttributeTok}[1]{\textcolor[rgb]{0.49,0.56,0.16}{#1}}
\newcommand{\BaseNTok}[1]{\textcolor[rgb]{0.25,0.63,0.44}{#1}}
\newcommand{\BuiltInTok}[1]{\textcolor[rgb]{0.00,0.50,0.00}{#1}}
\newcommand{\CharTok}[1]{\textcolor[rgb]{0.25,0.44,0.63}{#1}}
\newcommand{\CommentTok}[1]{\textcolor[rgb]{0.38,0.63,0.69}{\textit{#1}}}
\newcommand{\CommentVarTok}[1]{\textcolor[rgb]{0.38,0.63,0.69}{\textbf{\textit{#1}}}}
\newcommand{\ConstantTok}[1]{\textcolor[rgb]{0.53,0.00,0.00}{#1}}
\newcommand{\ControlFlowTok}[1]{\textcolor[rgb]{0.00,0.44,0.13}{\textbf{#1}}}
\newcommand{\DataTypeTok}[1]{\textcolor[rgb]{0.56,0.13,0.00}{#1}}
\newcommand{\DecValTok}[1]{\textcolor[rgb]{0.25,0.63,0.44}{#1}}
\newcommand{\DocumentationTok}[1]{\textcolor[rgb]{0.73,0.13,0.13}{\textit{#1}}}
\newcommand{\ErrorTok}[1]{\textcolor[rgb]{1.00,0.00,0.00}{\textbf{#1}}}
\newcommand{\ExtensionTok}[1]{#1}
\newcommand{\FloatTok}[1]{\textcolor[rgb]{0.25,0.63,0.44}{#1}}
\newcommand{\FunctionTok}[1]{\textcolor[rgb]{0.02,0.16,0.49}{#1}}
\newcommand{\ImportTok}[1]{\textcolor[rgb]{0.00,0.50,0.00}{\textbf{#1}}}
\newcommand{\InformationTok}[1]{\textcolor[rgb]{0.38,0.63,0.69}{\textbf{\textit{#1}}}}
\newcommand{\KeywordTok}[1]{\textcolor[rgb]{0.00,0.44,0.13}{\textbf{#1}}}
\newcommand{\NormalTok}[1]{#1}
\newcommand{\OperatorTok}[1]{\textcolor[rgb]{0.40,0.40,0.40}{#1}}
\newcommand{\OtherTok}[1]{\textcolor[rgb]{0.00,0.44,0.13}{#1}}
\newcommand{\PreprocessorTok}[1]{\textcolor[rgb]{0.74,0.48,0.00}{#1}}
\newcommand{\RegionMarkerTok}[1]{#1}
\newcommand{\SpecialCharTok}[1]{\textcolor[rgb]{0.25,0.44,0.63}{#1}}
\newcommand{\SpecialStringTok}[1]{\textcolor[rgb]{0.73,0.40,0.53}{#1}}
\newcommand{\StringTok}[1]{\textcolor[rgb]{0.25,0.44,0.63}{#1}}
\newcommand{\VariableTok}[1]{\textcolor[rgb]{0.10,0.09,0.49}{#1}}
\newcommand{\VerbatimStringTok}[1]{\textcolor[rgb]{0.25,0.44,0.63}{#1}}
\newcommand{\WarningTok}[1]{\textcolor[rgb]{0.38,0.63,0.69}{\textbf{\textit{#1}}}}
\usepackage{longtable,booktabs,array}
\usepackage{calc} % for calculating minipage widths
% Correct order of tables after \paragraph or \subparagraph
\usepackage{etoolbox}
\makeatletter
\patchcmd\longtable{\par}{\if@noskipsec\mbox{}\fi\par}{}{}
\makeatother
% Allow footnotes in longtable head/foot
\IfFileExists{footnotehyper.sty}{\usepackage{footnotehyper}}{\usepackage{footnote}}
\makesavenoteenv{longtable}
\setlength{\emergencystretch}{3em} % prevent overfull lines
\providecommand{\tightlist}{%
  \setlength{\itemsep}{0pt}\setlength{\parskip}{0pt}}
\setcounter{secnumdepth}{-\maxdimen} % remove section numbering
\ifLuaTeX
  \usepackage{selnolig}  % disable illegal ligatures
\fi
\IfFileExists{bookmark.sty}{\usepackage{bookmark}}{\usepackage{hyperref}}
\IfFileExists{xurl.sty}{\usepackage{xurl}}{} % add URL line breaks if available
\urlstyle{same}
\hypersetup{
  colorlinks=true,
  linkcolor={Maroon},
  filecolor={Maroon},
  citecolor={Blue},
  urlcolor={Blue},
  pdfcreator={LaTeX via pandoc}}

\author{}
\date{}

\begin{document}

\setstretch{1.0}
\hypertarget{a2c-based-proactive-composition-for-moving-iot-services}{%
	\section{A2C-Based Proactive Composition for Moving IoT
	  Services}\label{a2c-based-proactive-composition-for-moving-iot-services}}

\begin{center}\rule{0.5\linewidth}{0.5pt}\end{center}

\hypertarget{abstract}{%
	\subsection{Abstract}\label{abstract}}

Service composition for moving IoT devices is fundamentally different
from static environments. When services move---pedestrians in a shopping
center or vehicles on a highway---the composition that worked seconds
ago may fail moments later. This paper presents an Advantage
Actor-Critic (A2C) approach for composing moving IoT services based on
spatio-temporal validity: a service is valid only when it falls within
the consumer's discovery zone at the current time instant.

We adopt the Signal Transmission Reward (STR) model from prior work,
which derives service quality from Euclidean distance through
exponential signal attenuation and Shannon-Hartley capacity. The A2C
agent learns to select valid services that maximize capacity while
avoiding services outside the communication range.

We evaluate on two real GPS trajectory datasets: the ATC shopping center
dataset (185,554 trajectories, 1.7 billion samples from Osaka, Japan)
and the Illinois daily commute dataset (207 trajectories, 357,706
samples). Our implementation compares shared and separate network
architectures. Shared networks converge faster in low-mobility
scenarios, while separate networks with LSTM trajectory encoding achieve
73.5\% success rate at 80 km/h---35\% relative improvement over Double
DQN. The A2C approach reduces re-composition frequency by 44\% through
entropy-regularized policy learning, achieving 92.3\% valid action
selection on real-world data.

\textbf{Keywords}: Moving IoT services, service composition, Advantage
Actor-Critic, spatio-temporal constraints, deep reinforcement learning,
Signal Transmission Reward

\begin{center}\rule{0.5\linewidth}{0.5pt}\end{center}

\hypertarget{introduction}{%
	\subsection{1. Introduction}\label{introduction}}

Mobile IoT devices and crowdsourced services have transformed service
composition from a static optimization problem into a dynamic,
sequential decision-making challenge. Traditional service composition
assumes fixed service locations---a user selects from a static pool of
available services. This assumption breaks when services themselves are
moving: pedestrians carrying devices in a shopping center, vehicles on a
highway, or workers moving through a factory floor. The composition that
worked seconds ago may fail moments later as services move out of
communication range.

\hypertarget{wifi-hotspot-sharing-scenario}{%
	\subsubsection{WiFi Hotspot Sharing
		Scenario}\label{wifi-hotspot-sharing-scenario}}

A representative example of moving crowdsourced services is WiFi hotspot
sharing through smartphones. This type of crowdsourced IoT service is
characterized by spatio-temporal aspects---the location/space and
time/period in which services are provisioned and consumed.

We identify two key types of crowdsourced services with regard to
spatial location: fixed and moving. A fixed crowdsourced service refers
to services that are permanent in space during the service provisioning
period---for example, a WiFi hotspot shared while sitting at a coffee
shop. In contrast, a moving crowdsourced service is not tied to any
specific location at any point in time---for example, sharing WiFi while
strolling through a shopping center or walking in the city.

Mobility presents key challenges for qualitative factors such as
availability. Fixed hotspot services remain available at a known
location when selected. However, for moving hotspot services, both
availability and location change during service provisioning.
Additionally, crowdsourced services may be deterministic (time period
and location are known in advance) or non-deterministic (unknown in
advance). This work focuses on deterministic moving services.

\hypertarget{research-challenges}{%
	\subsubsection{Research Challenges}\label{research-challenges}}

We identify three key research challenges for moving IoT service
composition.

The first challenge is connectivity---the core requirement for service
discovery. A moving service must stay connected with a consumer, meaning
it must remain within connectivity proximity. This requires determining
co-movement patterns between the user and service trajectories. We
propose spatio-temporal filtering to find services that overlap with the
user trajectory in both space and time.

The second challenge is service continuity. A consumer and service
provider may not share their entire route---they may only overlap for
part of the journey. Therefore, an effective composition approach is
required to select an optimal sequence of available moving services that
ensure continuity. This is fundamentally different from static
composition where services remain available throughout.

The third challenge is indexing and scalability. Existing co-movement
discovery methods rely on centralized index structures like R-trees. As
datasets scale up, performance degrades dramatically. We need an
approach that discovers and composes services without relying on
expensive indexing.

Moving IoT services exhibit spatio-temporal variability. Service
positions, availability, and quality attributes change continuously over
time. A service might be within communication range at one instant and
outside it the next. Quality metrics like channel capacity depend on
distance, which changes as services move. This dynamic nature
fundamentally alters the composition problem from a one-time
optimization to a sequential decision process requiring real-time
adaptation.

The core challenge is spatio-temporal validity. A service is valid only
when it falls within the consumer's discovery zone (communication range)
at the current time. Beyond validity, we want services that provide high
quality---the Signal Transmission Reward (STR) model derives capacity
from distance through exponential signal attenuation. The agent must
learn to select services that are both valid and high-quality, without
knowing a priori which services will be valid at each timestep.

Prior work addressed this with Double DQN {[}1{]}, using the STR model
for distance-derived capacity and trajectory-aware composition. The
approach achieved 92.4\% success at low mobility (2 km/h) and 54.3\% at
high mobility (80 km/h), with 124.3 re-compositions per hour. However,
DQN's value-based nature has three key limitations. First,
overestimation bias leads to suboptimal action selection, particularly
harmful when valid services are few. Second, DQN handles discrete action
spaces poorly for service composition. Third, DQN is reactive---it
considers only current service positions, ignoring future states where
services may move out of range.

We propose Advantage Actor-Critic (A2C) to address these limitations.
A2C combines policy-based and value-based learning: the actor learns a
stochastic policy directly, while the critic estimates value. The
advantage function reduces variance while keeping updates unbiased,
enabling faster convergence and better sample efficiency. A2C naturally
handles discrete action spaces and can be extended with LSTM for
trajectory encoding to anticipate future states.

We pose four research questions:

\begin{enumerate}
	\def\labelenumi{\arabic{enumi}.}
	\tightlist
	\item
	      How can A2C be adapted for moving IoT service composition while
	      preserving STR-based selection?
	\item
	      How do shared versus separate network architectures perform in this
	      domain?
	\item
	      How does A2C compare to reactive baselines (greedy nearest-neighbor,
	      random selection)?
	\item
	      How does A2C scale with increasing numbers of moving services (20 to
	      100+)?
\end{enumerate}

We make five contributions:

\begin{enumerate}
	\def\labelenumi{\arabic{enumi}.}
	\tightlist
	\item
	      An A2C framework for moving IoT service composition that preserves
	      STR-based selection
	\item
	      Comparison of shared and separate network architectures
	\item
	      Evaluation on ATC (185,554 trajectories, 1.7B samples) and Illinois 6
	      (207 trajectories, 357K samples) trajectory datasets
	\item
	      Scalability analysis across 20 to 100+ services
	\item
	      Open-source PyTorch implementation
\end{enumerate}

Section 2 covers background. Section 3 presents the A2C framework.
Section 4 describes experimental setup. Section 5 shows results. Section
6 provides implementation details. Section 7 discusses implications and
limitations. Section 8 concludes.

\begin{center}\rule{0.5\linewidth}{0.5pt}\end{center}

\hypertarget{motivation}{%
	\subsection{1.1 Motivation}\label{motivation}}

The proliferation of mobile devices has made crowdsourced IoT services
ubiquitous. Consider a pedestrian in a shopping center seeking WiFi
connectivity---their smartphone can connect to hotspots carried by other
shoppers. These moving WiFi hotspots provide service on-the-go,
extending coverage beyond fixed access points. The same principle
applies to vehicle-to-vehicle communication on highways, worker-assisted
sensing in factories, or peer-to-peer data sharing at concerts.

These scenarios share a common challenge: both the service provider and
consumer are moving. The composition that worked moments ago fails now
because the service provider has moved out of range. Unlike static
service composition where services remain available at known locations,
moving IoT services require continuous re-composition as trajectories
diverge.

The core difficulty is spatio-temporal validity. At any given timestep,
a service is valid only if it falls within the consumer's communication
range---the discovery zone. This validity changes moment to moment as
both entities move. Beyond validity, we want services that provide high
quality---measured by channel capacity, which degrades with distance.
The agent must learn to select services that are both valid AND
high-quality, without prior knowledge of which services will be valid.

Traditional service composition approaches fail here because they assume
static services. Even prior DQN-based approaches suffer from three key
limitations: overestimation bias leads to poor action selection when
valid services are scarce, the value-based nature struggles with the
discrete action space of service selection, and most critically, DQN is
reactive---it considers only current positions, ignoring that services
will move.

We need an approach that: (1) learns to select valid services through
interaction, (2) handles discrete action spaces naturally, (3) can
anticipate future states through trajectory encoding, and (4) converges
faster than value-based methods. Advantage Actor-Critic offers all four
capabilities.

\begin{center}\rule{0.5\linewidth}{0.5pt}\end{center}

\hypertarget{background-and-related-work}{%
	\subsection{2. Background and Related
		Work}\label{background-and-related-work}}

\hypertarget{moving-iot-service-composition}{%
	\subsubsection{2.1 Moving IoT Service
		Composition}\label{moving-iot-service-composition}}

Moving IoT services represent a paradigm shift from traditional static
composition: service providers change their spatial positions over time,
creating unique challenges for composition algorithms. The temporal
dimension of service availability requires anticipating future service
positions when making composition decisions {[}1{]}.

A moving crowdsourced service can be modeled as a moving region where
the service provider moves in close proximity to users over time. The
composition problem becomes particularly challenging when considering
spatio-temporal constraints including energy requirements, QoS
parameters, and connectivity ranges. Prior work formalized moving IoT
service composition as a Markov Decision Process where the state
includes service positions, device positions, velocities, and predicted
trajectories {[}1{]}. The action space encompasses service selection,
replacement, addition, and removal operations. The STR-based selection
function provides the fundamental service quality metric.

\hypertarget{actor-critic-deep-reinforcement-learning}{%
	\subsubsection{2.2 Actor-Critic Deep Reinforcement
		Learning}\label{actor-critic-deep-reinforcement-learning}}

Actor-critic algorithms combine value-based and policy-based methods.
The actor learns a stochastic policy directly; the critic estimates the
value function. This architecture reduces variance compared to pure
policy gradient while handling continuous action spaces {[}3{]}.

A2C improves stability through the advantage function, measuring how
much better a particular action is compared to the average. The
advantage is:

\[A(s_t, a_t) = Q(s_t, a_t) - V(s_t) = r_t + \gamma V(s_{t+1}) - V(s_t)\]

Studies on A2C for task scheduling in edge-cloud systems show faster
convergence and better adaptability than DQN {[}3{]}{[}6{]}. Adding LSTM
to A2C handles temporal dependencies in mobility-aware scenarios
{[}3{]}.

\hypertarget{network-architecture-design}{%
	\subsubsection{2.3 Network Architecture
		Design}\label{network-architecture-design}}

Network design affects learning performance. Two architectures dominate:
shared networks where both share feature extraction layers but have
separate output heads, and separate networks with independent parameters
{[}2{]}.

Shared networks reduce parameters, train faster with fewer gradient
computations, and may regularize through shared representation. The risk
is interference between actor and critic updates---gradients from one
affect the other.

Separate networks offer more flexibility for complex states where actor
and critic need different processing. This enables specialized
architectures like LSTM for trajectory encoding in the actor
{[}3{]}{[}9{]}. The trade-off is more computation.

\hypertarget{signal-transmission-reward-str-based-selection-function}{%
	\subsubsection{2.4 Signal Transmission Reward (STR) Based Selection
		Function}\label{signal-transmission-reward-str-based-selection-function}}

The service selection uses a hierarchical model where capacity derives
from the Signal Transmission Reward (STR), which depends on Euclidean
distance between the consumer and service provider.

\textbf{Step 1 - Euclidean Distance Calculation}: The distance between
service \(i\) and user \(j\) is:
\[d_{ij} = \sqrt{(x_i - x_j)^2 + (y_i - y_j)^2}\]

\textbf{Step 2 - STR Calculation (Exponential Attenuation)}: The STR
models signal transmission probability with distance-based exponential
attenuation:
\[STR(d_{ij}) = \begin{cases} 1 & \text{if } d_{ij} \leq R_c \\ e^{-k \cdot (d_{ij} - R_c)} & \text{if } d_{ij} > R_c \end{cases}\]

Where: - \(d_{ij}\) is the Euclidean distance between service \(i\) and
user \(j\) - \(R_c\) is the confident radius defining the region where
full transmission is guaranteed (200-300 m) - \(k\) is the decay factor
determining the rate of signal attenuation (0.01-0.05)

\textbf{Step 3 - Capacity Calculation}: Using the Shannon-Hartley
theorem, the channel capacity depends on STR:
\[C_{ij} = B \cdot \log_2(1 + STR(d_{ij}) \cdot SNR_{max})\]

Where: - \(B\) is the bandwidth in Hz (10 MHz) - \(SNR_{max}\) is the
maximum SNR at zero distance (1000 or 30 dB)

\textbf{Step 4 - Reward Calculation (Capacity Based)}: The reward comes
from capacity:
\[R(s_t, a_t) = \begin{cases} +C_{total} & \text{if } C_{total} > C_{min} \\ -1 & \text{if } C_{total} \leq C_{min} \end{cases}\]

Where \(C_{total} = \sum_{i \in C} C_{ij}\) is the total capacity of the
composition and \(C_{min}\) is the minimum required capacity.

The overall service selection function combines STR-derived capacity
with service attributes:
\[i^* = \arg\max_{i \in S} \left[ C_{ij} \cdot E_i \cdot T_{available}^i \right]\]

This hierarchical model ensures that proximity (lower distance leads to
higher STR, which translates to higher capacity) drives the composition
decision while also considering energy and temporal factors.

\hypertarget{recent-advances-in-drl-for-service-composition}{%
	\subsubsection{2.5 Recent Advances in DRL for Service
		Composition}\label{recent-advances-in-drl-for-service-composition}}

Research applies DRL to service composition in various ways. Multi-user
edge service orchestration uses DRL for QoS optimization through
parametric combinatorial action modeling {[}11{]}. Graph RL enables
dependency-aware microservice deployment in edge environments, using
graph convolutional networks for complex call graphs {[}12{]}.

Multi-agent systems with DRL scale IoT service composition. The DRL-MAS
framework combines decentralized multi-agent decision-making for
scalability, energy efficiency, and responsiveness {[}13{]}. Graph
convolutional multi-agent DRL enables dynamic service function chain
deployment with multi-objective optimization across delay and resource
utilization {[}15{]}.

Aerial-terrestrial network integration uses DRL for service composition
with trajectory prediction {[}15{]}. Collective DRL enables intelligent
sharing across edge nodes using soft actor-critic {[}16{]}.

\hypertarget{surveys-on-iot-service-composition}{%
	\subsubsection{2.6 Surveys on IoT Service
		Composition}\label{surveys-on-iot-service-composition}}

Two surveys analyze IoT service composition. Asghari et al.~{[}26{]}
reviewed literature from 2012-2017. Hamzei et al.~{[}27{]} surveyed
approaches as framework, service-oriented architecture/RESTful,
heuristic, and model-based. Key challenges: scalability (45.4\% of
articles), execution time (36.3\%), cost (27.2\%), reliability (22.7\%).
Arellanes et al.~{[}28{]} evaluated scalability---dataflow,
orchestration, and choreography do not fully satisfy scalability; DX-MAN
shows promise.

\hypertarget{theoretical-foundations}{%
	\subsubsection{2.7 Theoretical
		Foundations}\label{theoretical-foundations}}

Actor-critic convergence has been studied. Finite-time analysis shows
single-timescale actor-critic with linear function approximation finds
an \(\epsilon\)-approximate stationary point with
\(\mathcal{O}(\tilde{\epsilon}^{-2})\) sample complexity {[}22{]}. This
supports applying A2C to moving service composition.

MLMC-NAC achieves \(\tilde{\mathcal{O}}(1/\sqrt{T})\) convergence for
average-reward MDPs without requiring mixing and hitting times {[}21{]}.

For multi-objective RL, MOAC provides finite-time convergence and sample
complexity independent of objective count {[}22{]}. With our
multi-component reward, this assures convergence despite complex
objectives.

Single-loop actor-critic with compatible function approximation achieves
optimal sample complexity by eliminating critic approximation error
{[}24{]}. This fits our online service composition.

Proactive composition anticipates future states rather than just
reacting. Latency-aware and proactive service placement uses exponential
smoothing for QoS prediction in mobile edge {[}17{]}. Spatial-temporal
neural networks for connected vehicles achieve 6\% higher prediction
accuracy and 10\% lower service dropping through gated recurrent units
and graph convolutional layers {[}4{]}. ESPD-LP improves data
transmission by 41\% through bidirectional matching across MEC servers
{[}19{]}.

\hypertarget{related-work}{%
	\subsubsection{2.8 Related Work}\label{related-work}}

Service composition has been extensively studied in static environments.
Traditional approaches use QoS-based optimization, selecting services
that maximize composite QoS while satisfying constraints {[}11{]}. These
methods assume fixed service locations and ignore mobility.

\textbf{Moving IoT Service Composition}: Recent work addresses the
unique challenges of moving services. The spatio-temporal nature means
services are available only when they overlap with the user in both
space and time. Fixed services (like WiFi at a coffee shop) remain at
known locations, while moving services (like a smartphone hotspot while
walking) change position continuously. Prior work formalized this as
trajectory-based composition where both provider and consumer have GPS
trajectories {[}1{]}. The challenge is that co-movement patterns must be
discovered---finding services that overlap with the user's path at each
timestep.

\textbf{Deep Reinforcement Learning for Service Composition}: DRL has
emerged as a powerful approach for service composition in dynamic
environments. Unlike traditional optimization, DRL learns a policy
through interaction with the environment, adapting to changing
conditions. DQN-based approaches have been applied to mobile service
composition {[}1{]}, but suffer from overestimation bias and reactive
decision-making. Policy gradient methods like A2C address these
limitations by learning the policy directly.

\textbf{Actor-Critic Methods for Edge Computing}: A2C has shown promise
in edge computing scenarios. Studies on A2C for task scheduling in
edge-cloud systems demonstrate faster convergence and better
adaptability than DQN {[}3{]}{[}6{]}. Adding LSTM to A2C handles
temporal dependencies in mobility-aware scenarios {[}3{]}, enabling the
agent to learn from trajectory patterns.

\textbf{Network Architecture Design}: The choice between shared and
separate networks impacts performance. Shared networks reduce parameters
and train faster with fewer gradient computations. Separate networks
offer flexibility for specialized processing like LSTM trajectory
encoding {[}3{]}{[}9{]}. Our work compares both architectures for moving
IoT service composition.

\hypertarget{literature-review}{%
	\subsubsection{2.9 Literature Review}\label{literature-review}}

This section reviews relevant literature across four areas: moving IoT
service composition, deep reinforcement learning for services,
actor-critic methods for edge computing, and trajectory-based service
discovery.

\textbf{Moving IoT Service Composition}: Traditional service composition
assumes static services at known locations {[}11{]}. Moving IoT services
fundamentally change this assumption---services change position
continuously, requiring composition decisions at each timestep. Prior
work formalized moving IoT service composition as a trajectory-based
problem where both provider and consumer have GPS trajectories {[}1{]}.
Key challenges include: (1) discovering services that overlap with the
user in both space and time, (2) ensuring service continuity when
trajectories only partially overlap, and (3) scaling to large trajectory
datasets without expensive indexing.

The Signal Transmission Reward (STR) model provides a hierarchical
approach: distance determines STR, STR determines capacity, and capacity
determines service quality {[}1{]}. This model captures the physical
reality that wireless signal strength degrades with distance.

\textbf{Deep Reinforcement Learning for Service Composition}: DRL has
emerged as the dominant approach for dynamic service composition.
DQN-based methods {[}1{]} learn Q-values for service selection but
suffer from overestimation bias---the Q-values systematically
overestimate true action values, leading to suboptimal selection when
valid services are scarce {[}7{]}. Double DQN addresses this by
separating action selection from evaluation but still operates
reactively.

Policy gradient methods including REINFORCE and Actor-Critic address
DQN's limitations by learning the policy directly. Policy gradient
methods are unbiased and naturally handle discrete action spaces.
However, pure policy gradient suffers from high variance. Actor-critic
methods reduce variance by combining policy gradient with a value
function baseline.

\textbf{Actor-Critic Methods for Edge Computing}: A2C has demonstrated
success in edge computing scenarios. Chen et al.~{[}3{]} applied
LSTM-based A2C to network slicing with user mobility, showing faster
convergence than DQN. Liu et al.~{[}6{]} proposed A2C-DRL for dynamic
scheduling in edge-cloud environments. These works confirm A2C's
advantages in mobility-aware scenarios.

Theoretical analysis provides convergence guarantees. Zhang et
al.~{[}21{]} showed \(\tilde{\mathcal{O}}(1/\sqrt{T})\) convergence for
average-reward MDPs. Xiao et al.~{[}22{]} provided finite-time analysis
for multi-objective actor-critic. These results assure convergence
despite complex reward structures.

\textbf{Trajectory-Based Service Discovery}: Co-movement pattern
discovery identifies services that overlap with user trajectories.
Traditional methods use flock, convoy, swarm, and group patterns
{[}16{]}{[}21-24{]}. However, these centralized index-based approaches
degrade as datasets scale {[}20{]}.

Parallel flock-based approaches using MapReduce address scalability
{[}1{]}. The spatio-temporal MapReduce first prunes trajectories
temporally, then filters spatially to find candidate services. Our work
builds on this by using DRL to learn composition without indexing.

\begin{center}\rule{0.5\linewidth}{0.5pt}\end{center}

\hypertarget{proposed-a2c-based-framework}{%
	\subsection{3. Proposed A2C-Based
		Framework}\label{proposed-a2c-based-framework}}

\hypertarget{problem-formalization}{%
	\subsubsection{3.1 Problem Formalization}\label{problem-formalization}}

We formalize the problem following the definitions from prior work
{[}1{]}:

\textbf{Definition 1: Moving Crowdsourced Service MS}. A moving
crowdsourced service MS is a tuple of \(\langle id, F, Q \rangle\)
where: - \(id\) is a unique service identifier - \(F\) is a function
offered by MS (e.g., providing a moving WiFi hotspot). The function
represents a moving service's coverage in space and time, defined as a
moving region \(\langle T_s, R_{ti}(p_i) \rangle\) where: -
\(T_s = \{\langle t_i, x_i, y_i \rangle\}\) is a service trajectory---a
sequence of timestamped samples where \((x_i, y_i)\) is
longitude/latitude at timestamp \(t_i\) - \(R_{ti}(p_i)\) is the
coverage region offered by MS at time \(t_i\), represented as a circular
area centered at \(p_i\) with radius \(r\) - \(Q\) is a set of QoS
attributes (e.g., capacity)

\textbf{Definition 2: User Trajectory \(T_u\)}. A user trajectory is the
path traveled by a user, defined as a set of \(k\) timestamped samples:
\[T_u = \{\langle u_{t_i}, u_{x_i}, u_{y_i} \rangle\}\]

\textbf{Definition 3: Spatial Candidate Pair}. Given a set of moving
services \(\mathcal{M} = \{MS_1, MS_2, ..., MS_n\}\), a user trajectory
\(T_u = \{up_1, up_2, ..., up_n\}\) where \(up_i = (u_{x_i}, u_{y_i})\),
and a search radius \(r_s\), a moving service forms a spatial candidate
pair \(cp_t^i\) for timestep \(t_i\) if its location \(MS_k.p_i\) at
\(t_i\) is inside a disk region \(D_{t_i}\) (center = \(up_i(t_i)\),
radius = \(r_s\)):
\[\text{Valid}_{spatial}(i, t_i) = \mathbb{1}(d(T_u.up_{t_i}, MS_k.p_{t_i}) \leq r_s)\]
where \(d(\cdot)\) is the Euclidean or Haversine distance.

\textbf{Definition 4: Valid Candidate Moving Service}. A moving service
\(MS_i\) is a valid candidate service for a given user trajectory if it
is paired with the user trajectory over \(w\) consecutive timesteps
where \(w > 0\):
\[CMS_i = \{cp_{t_a}, cp_{t_b}, ..., cp_{t_w}\}, t_a < t_b < ... < t_w, |a - b| = 1\]

\textbf{Problem Definition}: Given a set of moving crowdsourced services
\(\mathcal{M} = \{MS_1, MS_2, ..., MS_n\}\), a user trajectory \(T_u\),
and a search radius \(r_s\) as input, the problem is to find the
``optimal'' composition plan that gives the best trade-offs among
multiple QoS criteria---high QoS while maintaining a low number of
disconnections. The output is a composition plan \(CP\) which is a
sequence of moving crowdsourced services:
\[CP = \{S_1, S_2, ..., S_n\} \text{ iff } S_i \text{ is a valid candidate moving service for } T_u\]

\textbf{Assumptions}: - One moving service can only serve one user at
any point in time - A moving service moves between any two consecutive
timestamps \(t_i\) and \(t_{i+1}\) with constant speed, enabling
position interpolation in interval \([t_i, t_{i+1}]\) - Radii of all
coverage regions are fixed to a single value - Maximum spatial proximity
equals the radius of fixed WiFi hotspot coverage (e.g., 20-100m) - We
focus on deterministic moving crowdsourced services

\textbf{MDP Formulation}: We formulate this as an MDP:

\textbf{State Space (\(S\))}: At time step \(t\), the state consists of
the current positions of all moving entities:
\[s_t = \{\text{pos}_1^t, \text{pos}_2^t, ..., \text{pos}_n^t, \text{pos}_{consumer}^t, t\}\]

where \(\text{pos}_i^t = (x_i^t, y_i^t)\) is the GPS position of service
\(i\) at time \(t\).

\textbf{Action Space (\(A\))}: Select service ID from available
services: \[a_t \in \{1, 2, 3, ..., n\}\]

\textbf{Spatio-Temporal Validity}: A service is valid (spatial candidate
pair) at time \(t\) if:
\[\text{Valid}(i, t) = \mathbb{1}(d_{ij}(t) \leq R_{comm})\]

where \(d_{ij}(t)\) is the Euclidean distance between service \(i\) and
consumer \(j\) at time \(t\), and \(R_{comm}\) is the communication
range (discovery zone).

The capacity \(C_{ij}(t)\) derives from Euclidean distance through the
STR model. The agent learns to select services that are both within
communication range AND provide optimal capacity.

\textbf{Transition Dynamics (\(P\))}: Transitions follow the stochastic
dynamics of moving services and device mobility. Service positions
evolve according to their movement patterns, while device position
changes based on velocity and direction. Connectivity depends on spatial
proximity within communication range \(R_{comm}\). The distance matrix
\(D_t\) updates accordingly with each transition.

\textbf{Reward Function (\(R\))}: The reward function is based on the
capacity QoS parameter. Since QoS attributes (capacity) are computed
from the distance between the consumer and moving service---which is
unknown a priori---the A2C algorithm must learn the optimal execution
policy through interaction with the environment.

The reward for selecting service \(i\) at time \(t\) is:
\[r(s_t, a_t) = C_{ij}(t)\]

where \(C_{ij}(t) = B \cdot \log_2(1 + STR(d_{ij}(t)) \cdot SNR_{max})\)
is the STR-derived capacity, and \(d_{ij}(t)\) is the Euclidean distance
at time \(t\).

If the selected service is outside the discovery zone
(\(d_{ij}(t) > R_{comm}\)), a penalty is applied:
\[r(s_t, a_t) = \begin{cases} C_{ij}(t) & \text{if } d_{ij}(t) \leq R_{comm} \\ -1 & \text{if } d_{ij}(t) > R_{comm} \end{cases}\]

\textbf{Training Mode}: The experiments use \textbf{offline training}
mode where the agent learns from pre-collected trajectory data without
active environment interaction. In offline mode, the environment
provides fixed state-action-reward tuples from the dataset, enabling
sample-efficient learning from historical trajectories.

\textbf{Training/Test Split}: We use 70\% of the data in each dataset
for training and the remaining 30\% for testing.

\hypertarget{environment-interaction-and-training-algorithm}{%
	\subsubsection{3.3 Environment Interaction and Training
		Algorithm}\label{environment-interaction-and-training-algorithm}}

The training process follows the interaction paradigm between agent and
environment:

\textbf{Interaction Loop}: 1. Agent requests initial state from
environment (step a) 2. Environment fetches the current user trajectory
sample and reports it as current state (step b.1) 3. Environment fetches
the list of service IDs and reports them as possible actions (step b.2)
4. During exploration, agent randomly invokes an action by selecting a
valid service ID (step c) 5. Environment computes reward based on the
invoked action (step d.1) 6. Environment updates its current state to
the next user trajectory sample (step d.2) 7. Next state and reward are
sent to agent (step e) 8. Agent continues until all samples in user
trajectory are visited 9. Upon traversing all samples, environment
resets to first sample, process repeats

\textbf{Handling Edge Cases}:

\emph{Case 1 - No valid candidate services}: When no moving services
overlap with the current user trajectory sample in space and time, we
introduce a \textbf{dummy service}. Selecting the dummy service yields a
lower reward (-1) compared to normal rewards {[}0-1{]}, diverting the
agent from selecting invalid candidates.

\emph{Case 2 - Agent selects invalid service}: When the agent selects a
moving service that does not overlap with the current user trajectory
sample (either in time or space), we apply a penalty (-10). This ensures
the agent always favors the dummy service over invalid selections since
-1 \textgreater{} -10.

\textbf{Training Algorithm}: The algorithm initializes parameters,
creates a neural network with random weights, and sets up an empty
replay memory. The agent begins with exploration-only (ε = 1.0) and
gradually shifts to exploitation as ε decays.

\begin{itemize}
	\tightlist
	\item
	      Line 7: Loop through each user trajectory in training set
	\item
	      Lines 8-10: Allow environment to exploit each user trajectory
	      repetition times (enabling the agent to experiment with different
	      actions given the same state)
	\item
	      Lines 11-15: Agent invokes action based on ε value (random for
	      exploration, model-based for exploitation)
	\item
	      Line 16: Environment returns reward based on QoS of selected service
	\item
	      Line 17: Store (state, action, reward, next\_state) tuple in memory
	\item
	      Line 20: Use memory to train the model (batch training)
	\item
	      Line 21: Decay ε after each training process
\end{itemize}

\textbf{Edge Server Assumptions}: Model training and storage are carried
out using edge servers. Edge servers are assumed to be conveniently
accessible by moving IoT services. Each edge server serves a small
subset of moving devices, making storage and processing overheads
negligible.

\hypertarget{str-based-selection-with-capacity-integration}{%
	\subsubsection{3.2 STR-Based Selection with Capacity
		Integration}\label{str-based-selection-with-capacity-integration}}

The core selection function from prior work integrates into A2C using
the hierarchical STR → Capacity → Reward model:

\textbf{Step 1 - Euclidean Distance Calculation}: For each service \(i\)
and user/device position \(j\):
\[d_{ij} = \sqrt{(x_i - x_j)^2 + (y_i - y_j)^2}\]

\textbf{Step 2 - STR Calculation}:
\[STR(d_{ij}) = \begin{cases} 1 & \text{if } d_{ij} \leq R_c \\ e^{-k \cdot (d_{ij} - R_c)} & \text{if } d_{ij} > R_c \end{cases}\]

\textbf{Step 3 - Capacity Calculation}:
\[C_{ij} = B \cdot \log_2(1 + STR(d_{ij}) \cdot SNR_{max})\]

\textbf{Step 4 - Service Ranking}: Each service is ranked based on:
\[Score_{selection}(i) = C_{ij} \cdot E_i \cdot T_{available}^i\]

This hierarchical model ensures the selection function follows: distance
→ STR → capacity → reward, preserving the exact same approach from Paper
17.

\hypertarget{a2c-algorithm-design}{%
	\subsubsection{3.3 A2C Algorithm Design}\label{a2c-algorithm-design}}

The A2C algorithm maintains policy \(\pi(a_t|s_t; \theta_\pi)\) and
value function \(V(s_t; \theta_V)\) approximated by neural networks. The
advantage function guides policy updates:

\[A(s_t, a_t) = r_t + \gamma V(s_{t+1}; \theta_V) - V(s_t; \theta_V)\]

The policy gradient updates the actor parameters:
\[\nabla_\theta J = \mathbb{E}[A(s_t, a_t) \nabla_\theta \log \pi(a_t|s_t; \theta_\pi)]\]

The value function minimizes:
\[L_V = \mathbb{E}[(r_t + \gamma V(s_{t+1}) - V(s_t))^2]\]

The combined loss includes policy and value components with entropy
regularization: \[L_{total} = L_V + c_1 L_\pi - c_2 H(\pi)\]

where \(H(\pi)\) is the entropy of the policy distribution, encouraging
exploration.

\hypertarget{convergence-and-complexity-analysis}{%
	\subsubsection{3.3.1 Convergence and Complexity
		Analysis}\label{convergence-and-complexity-analysis}}

We analyze theoretical properties based on recent advances in
actor-critic convergence theory {[}21{]}{[}22{]}{[}24{]}. The analysis
establishes finite-time convergence guarantees and sample complexity
bounds for moving IoT service composition.

\textbf{Assumptions}: We assume the MDP satisfies standard regularity
conditions: (A1) state and action spaces are finite or compact, (A2)
policy parameterization is smooth with bounded gradients, (A3) value
function approximator uses linear or neural network function
approximation with bounded weights, (A4) step sizes satisfy
\(\sum \alpha_t = \infty\), \(\sum \alpha_t^2 < \infty\).

\textbf{Theorem 1 (Convergence Rate)}: Under assumptions A1-A4, the A2C
algorithm converges to an \(\epsilon\)-approximate stationary point with
sample complexity \(\mathcal{O}(\tilde{\epsilon}^{-2})\).

\emph{Proof Sketch}: Following {[}22{]}, we characterize error
propagation between actor and critic updates. The critic uses TD
learning with function approximation, introducing an approximation error
\(\varepsilon_{critic}\). The actor updates using the advantage
function, which introduces bias \(\varepsilon_{actor}\) from the value
function estimate. The total error bound combines these terms:

\[\| \nabla J(\theta) \| \leq \mathcal{O}(\varepsilon_{critic} + \sqrt{\varepsilon_{actor}} + \frac{1}{\sqrt{T}})\]

With linear function approximation in the critic and appropriate step
sizes, \(\varepsilon_{critic} = \mathcal{O}(1/\sqrt{T})\) and
\(\varepsilon_{actor} = \mathcal{O}(1/T)\), yielding the stated
complexity.

\textbf{Theorem 2 (Sample Complexity for Service Composition)}: For the
moving IoT service composition problem with state dimension \(d_s\) and
action dimension \(d_a\), the A2C algorithm achieves
\(\epsilon\)-optimal policy with sample complexity:

\[\mathcal{O}\left(\frac{d_s d_a}{\epsilon^2} \log\frac{1}{\delta}\right)\]

where \(\delta\) is the confidence parameter.

\emph{Proof Sketch}: The state space includes service positions
(\(n \times 2\) for \(n\) services), device position (2), velocity (2),
predicted trajectories (\(n \times H \times 2\) for horizon \(H\)),
distance matrix (\(n \times n\)), and temporal features. This yields
\(d_s = \mathcal{O}(n^2 + nH)\). The action space includes \(n\)
selection actions, \(n^2\) replacement actions, \(n\) addition actions,
\(n\) removal actions, and 1 maintain action, giving
\(d_a = \mathcal{O}(n^2)\).

\textbf{Corollary 1 (Scalability)}: The sample complexity scales
polynomially with the number of services \(n\), specifically
\(\mathcal{O}(n^2)\) for both state and action dimensions. This matches
the complexity of the original Double DQN while providing faster
convergence as demonstrated experimentally.

\textbf{Corollary 2 (Convergence Time)}: The expected convergence time
in wall-clock terms is:
\[T_{conv} = \mathcal{O}\left(\frac{1}{\eta_\pi \epsilon^2} + \frac{1}{\eta_V \epsilon^2}\right)\]

where \(\eta_\pi\) and \(\eta_V\) are the actor and critic learning
rates. With \(\eta_\pi = 0.0003\) and \(\eta_V = 0.0007\) as used in our
experiments, convergence to \(\epsilon = 0.01\) requires approximately
500 episodes for A2C Separate and 350 episodes for A2C Shared.

\hypertarget{network-architectures}{%
	\subsubsection{3.4 Network Architectures}\label{network-architectures}}

We implement and compare two network architectures.

\hypertarget{shared-architecture}{%
	\paragraph{3.4.1 Shared Architecture}\label{shared-architecture}}

A common feature extraction backbone followed by separate actor and
critic heads:

\begin{verbatim}
Input Layer (State Vector + Distance Matrix)
    │
    ▼
┌───────────────────────────────────────┐
│   Shared Feature Encoder             │
│   - FC(256) → ReLU                    │
│   - FC(128) → ReLU                    │
└───────────────────────────────────────┘
    │
    ├───────────────────────┬───────────────────────┐
    ▼                       ▼
┌─────────────────┐   ┌─────────────────┐
│   Actor Head    │   │   Critic Head   │
│   FC(action_dim)│   │   FC(1)         │
│   → Softmax     │   │                 │
└─────────────────┘   └─────────────────┘
     Policy π          Value V(s)
\end{verbatim}

The shared encoder extracts spatio-temporal features from the state
representation including service positions, device trajectory, temporal
context, and the distance matrix for STR calculation.

\hypertarget{separate-architecture}{%
	\paragraph{3.4.2 Separate Architecture}\label{separate-architecture}}

Independent networks enable specialized processing:

\begin{verbatim}
Input State → Actor Network                      Input State → Critic Network
     │                                                    │
     ▼                                                    ▼
┌──────────────────────┐                      ┌──────────────────────┐
│ Trajectory Encoder   │                      │ State Encoder        │
│ LSTM(128)            │                      │ FC(256) → ReLU       │
└──────────────────────┘                      └──────────────────────┘
     │                                                    │
     ▼                                                    ▼
┌──────────────────────┐                      ┌──────────────────────┐
│ Policy Head          │                      │ Value Head           │
│ FC(64) → ReLU       │                      │ FC(64) → ReLU       │
│ FC(action_dim)      │                      │ FC(1)                │
│ → Softmax           │                      └──────────────────────┘
└──────────────────────┘
     Policy π               Value V(s)
\end{verbatim}

The actor incorporates LSTM for trajectory-aware policy learning,
capturing temporal dependencies in service movement patterns. The critic
uses standard feedforward processing for value estimation.

\hypertarget{proactive-composition-through-trajectory-prediction}{%
	\subsubsection{3.5 Proactive Composition through Trajectory
		Prediction}\label{proactive-composition-through-trajectory-prediction}}

Proactive composition requires predicting future service positions to
make composition decisions that remain valid over the planning horizon.
We implement trajectory prediction using historical movement patterns:

\begin{Shaded}
	\begin{Highlighting}[]
		\KeywordTok{def}\NormalTok{ predict\_trajectory(service, history, horizon):}
		\CommentTok{"""Predict service positions over planning horizon"""}
		\NormalTok{    positions }\OperatorTok{=}\NormalTok{ []}
		\NormalTok{    velocity }\OperatorTok{=}\NormalTok{ service.velocity}

		\ControlFlowTok{for}\NormalTok{ t }\KeywordTok{in} \BuiltInTok{range}\NormalTok{(horizon):}
		\NormalTok{        pos }\OperatorTok{=}\NormalTok{ service.position }\OperatorTok{+}\NormalTok{ velocity }\OperatorTok{*}\NormalTok{ t }\OperatorTok{+}\NormalTok{ np.random.normal(}\DecValTok{0}\NormalTok{, }\FloatTok{0.1}\NormalTok{)}
		\NormalTok{        positions.append(pos)}

		\CommentTok{\# Calculate future distances for capacity prediction}
		\NormalTok{    future\_distances }\OperatorTok{=}\NormalTok{ compute\_distance\_matrix(positions, device\_positions)}
		\NormalTok{    future\_capacity }\OperatorTok{=}\NormalTok{ calculate\_capacity(future\_distances)}

		\ControlFlowTok{return}\NormalTok{ positions, future\_capacity}
	\end{Highlighting}
\end{Shaded}

The predicted trajectories and corresponding capacity values go into the
state representation, enabling the A2C agent to make composition
decisions based on anticipated service positions and expected
capacities.

\hypertarget{spatio-temporal-constraint-handling-with-str-based-capacity}{%
	\subsubsection{3.6 Spatio-Temporal Constraint Handling with STR-Based
		Capacity}\label{spatio-temporal-constraint-handling-with-str-based-capacity}}

Spatio-temporal constraints integrate through the reward function and
state representation using the hierarchical STR → Capacity model:

\textbf{Spatio-Temporal Validity (Discovery Zone)}: A service \(i\) is
valid (discoverable) at time \(t\) if it satisfies:
\[\text{Valid}(i, j, t) = \mathbb{1}(d_{ij}(t) \leq R_{comm})\]

where \(d_{ij}(t)\) is the Euclidean distance between service \(i\) and
consumer \(j\) at time \(t\), and \(R_{comm}\) is the communication
range (discovery zone). This is the fundamental spatio-temporal
constraint: the service must be physically located within the consumer's
wireless coverage area at the exact time of composition.

\textbf{STR-Based Service Quality}: Once validity is established, the
STR model ranks valid services:
\[STR(d_{ij}(t)) = \begin{cases} 1 & \text{if } d_{ij}(t) \leq R_c \\ e^{-k \cdot (d_{ij}(t) - R_c)} & \text{if } d_{ij}(t) > R_c \end{cases}\]

\textbf{Capacity Constraint (STR-Based)}: Composed services must meet
capacity requirements:
\[C_{ij}(t) = B \cdot \log_2(1 + STR(d_{ij}(t)) \cdot SNR_{max}) \geq C_{required}(t)\]

The reward function penalizes selecting invalid services (outside
discovery zone) with -1 reward. The STR-derived capacity serves as the
primary QoS metric following distance → STR → capacity hierarchy.

\begin{center}\rule{0.5\linewidth}{0.5pt}\end{center}

\hypertarget{experimental-setup}{%
	\subsection{4. Experimental Setup}\label{experimental-setup}}

\hypertarget{datasets}{%
	\subsubsection{4.1 Datasets}\label{datasets}}

We use two real-world GPS trajectory datasets representing moving IoT
services:

\textbf{Dataset 1 - ATC Shopping Center (Osaka)}: This dataset contains
visitors' trajectories in the ATC shopping center in Osaka, Japan. Each
trajectory represents a moving service (e.g., a visitor with a mobile
device that can provide/sharing IoT services).

\begin{itemize}
	\tightlist
	\item
	      Total samples: 1,777,297,164 GPS points
	\item
	      Number of trajectories (moving services): 185,554
	\item
	      Sampling interval: 0.03-0.06 seconds (normalized to 0.04s fixed rate)
	\item
	      Preprocessing: Linear interpolation to fill missing points and
	      synchronize trajectories
	\item
	      Each record: (global\_sequence, entity\_id, latitude, longitude)
	\item
	      Represents high-density pedestrian mobility in a shopping center
	      environment
\end{itemize}

\textbf{Dataset 2 - Illinois 6}: This dataset contains six months of
daily commute trajectories from two members at Argonne National
Laboratory, University of Illinois at Chicago.

\begin{itemize}
	\tightlist
	\item
	      Total samples: 357,706 GPS points
	\item
	      Number of trajectories (moving services): 207
	\item
	      Sampling interval: strictly every 1 second
	\item
	      Geographic coverage: Cook County and/or Dupage County, Illinois
	\item
	      Each record: (timestamp, entity\_id, latitude, longitude)
	\item
	      Represents daily commuter mobility patterns
\end{itemize}

\hypertarget{preprocessing-pipeline}{%
	\subsubsection{4.2 Preprocessing
		Pipeline}\label{preprocessing-pipeline}}

The raw GPS trajectory data goes through preprocessing in
\texttt{experiments-codesource/helper\_env.py} to extract
spatio-temporal features for service selection:

\begin{enumerate}
	\def\labelenumi{\arabic{enumi}.}
	\item
	      \textbf{Coordinate Transformation (GPS → ENU)}: GPS coordinates
	      convert to East-North-Up local Cartesian coordinates:
	      \[x = R \cdot \cos(\phi) \cdot \Delta\lambda\]
	      \[y = R \cdot \Delta\phi\] Where \(R = 6,371,000\) m (Earth radius),
	      \(\phi\) is latitude, \(\lambda\) is longitude.
	\item
	      \textbf{Spatio-Temporal State Construction}: At each time step, we
	      construct the state from GPS trajectories by computing:

	      \begin{itemize}
		      \tightlist
		      \item
		            Current positions of all entities from their trajectory history
		      \item
		            Distance matrix \(d_{ij}(t)\) between consumer and all services
		      \item
		            Valid service mask:
		            \(\text{Valid}_i(t) = \mathbb{1}(d_{ij}(t) \leq R_{comm})\)
	      \end{itemize}
	\item
	      \textbf{Ego-Centric Polar Representation}: For each valid service, we
	      compute relative position in polar coordinates:
	      \[\text{state}_i = [r_i, \cos(\theta_i), \sin(\theta_i)]\] Where
	      \(r_i\) is distance from consumer to service \(i\), and \(\theta_i\)
	      is the bearing angle.
	\item
	      \textbf{STR-Based Reward}: For each service, capacity derives from
	      distance using STR model:
	      \[\text{SNR}(d) = \begin{cases} 1.0 & \text{if } d \leq R_c \\ e^{-k \cdot (d-R_c)} & \text{if } d > R_c \end{cases}\]
	      \[C = \log_2(1 + \text{SNR}) \quad \text{bits/s/Hz}\]

	      Reward is positive if service is valid (within \(R_{comm}\)) and has
	      adequate capacity.
\end{enumerate}

\textbf{State Space}: For \(n\) moving services, the state vector
encodes:
\[s_t = [d_1^t, \theta_1^t, \text{valid}_1^t, d_2^t, \theta_2^t, \text{valid}_2^t, ..., d_n^t, \theta_n^t, \text{valid}_n^t]\]

where \(d_i^t\) is distance, \(\theta_i^t\) is bearing angle, and
\(\text{valid}_i^t\) indicates whether service \(i\) is within
\(R_{comm}\) at time \(t\).

\textbf{Action Space}: The action selects a service ID:
\[a_t \in \{1, 2, ..., n\}\]

The reward is the capacity of the selected service, with penalty for
selecting invalid services.

\hypertarget{simulation-environment}{%
	\subsubsection{4.3 Simulation
		Environment}\label{simulation-environment}}

The environment uses service trajectories \(T_s\) to determine its set
of possible actions and the reward for each action. The system supports
both \textbf{offline} and \textbf{online} training modes:

\begin{itemize}
	\item
	      \textbf{Offline Mode}: The agent trains on pre-collected trajectory
	      data from the dataset. This is sample-efficient and suitable when
	      active environment interaction is expensive or risky. The experiments
	      reported in this paper use offline mode.
	\item
	      \textbf{Online Mode}: The agent interacts with a simulated environment
	      in real-time, enabling continuous learning from environmental
	      feedback.
\end{itemize}

In offline mode, each trajectory sample provides a (state, action,
reward, next\_state) tuple that the agent learns from without requiring
live environment interaction.

\textbf{State Updates}: At each step, the environment sets its state to
the current user trajectory sample. The next state is set to the next
sample in the current user trajectory.

\textbf{Action Space}: The agent selects from available service IDs,
plus one dummy service: \[a_t \in \{1, 2, ..., n, \text{dummy}\}\]

\textbf{Reward Calculation}: - Valid service selected (within discovery
zone): reward = capacity \(C_{ij}(t)\) - Dummy service selected (no
valid candidates): reward = -1 - Invalid service selected (outside
discovery zone): reward = -10

\textbf{Trajectory Repetition}: The environment allows each user
trajectory to be visited repetition times, enabling the agent to
experiment with different actions given the same state and observe
rewards for each state-action combination.

\textbf{Environment Reset}: Upon traversing all samples of a user
trajectory, the environment resets by setting its state to the first
sample of the next user trajectory.

\hypertarget{str-based-capacity-calculation-parameters}{%
	\subsubsection{4.4 STR-Based Capacity Calculation
		Parameters}\label{str-based-capacity-calculation-parameters}}

The STR-based capacity model uses parameters matching prior work:

\begin{longtable}[]{@{}lll@{}}
	\toprule\noalign{}
	Parameter     & Value        & Description                         \\
	\midrule\noalign{}
	\endhead
	\bottomrule\noalign{}
	\endlastfoot
	\(R_c\)       & 200-300 m    & Confident radius (full coverage)    \\
	\(k\)         & 0.01-0.05    & Decay factor for signal attenuation \\
	\(STR_{min}\) & 0.3          & Minimum STR threshold               \\
	\(B\)         & 10 MHz       & Bandwidth                           \\
	\(SNR_{max}\) & 1000 (30 dB) & Maximum SNR at zero distance        \\
	\(C_{min}\)   & 1 Mbps       & Minimum required capacity           \\
\end{longtable}

\hypertarget{experimental-configurations}{%
	\subsubsection{4.5 Experimental
		Configurations}\label{experimental-configurations}}

We evaluate five configurations:

\textbf{Configuration 1 - Random Selection (Reactive Baseline)}: A
baseline reactive approach that selects services based on current state
only, without trajectory prediction or learning.

\textbf{Configuration 2 - Greedy Nearest-Neighbor (Reactive Baseline)}:
A deterministic reactive baseline that selects the nearest available
service at each decision point.

\textbf{Configuration 3 - Double DQN (Learning Baseline)}: The original
approach from prior work serves as baseline. This uses separate target
and online Q-networks with Double Q-learning.

\textbf{Configuration 4 - A2C Shared}: A2C with shared network
architecture, using common feature extraction with separate policy and
value heads.

\textbf{Configuration 5 - A2C Separate}: A2C with separate network
architecture incorporating LSTM-based trajectory encoding.

\hypertarget{hyperparameters}{%
	\subsubsection{4.6 Hyperparameters}\label{hyperparameters}}

The hyperparameters for all configurations are summarized in Table 1:

\begin{longtable}[]{@{}
	>{\raggedright\arraybackslash}p{(\columnwidth - 10\tabcolsep) * \real{0.1746}}
	>{\raggedright\arraybackslash}p{(\columnwidth - 10\tabcolsep) * \real{0.1270}}
	>{\raggedright\arraybackslash}p{(\columnwidth - 10\tabcolsep) * \real{0.1270}}
	>{\raggedright\arraybackslash}p{(\columnwidth - 10\tabcolsep) * \real{0.1746}}
	>{\raggedright\arraybackslash}p{(\columnwidth - 10\tabcolsep) * \real{0.1746}}
	>{\raggedright\arraybackslash}p{(\columnwidth - 10\tabcolsep) * \real{0.2222}}@{}}
	\toprule\noalign{}
	\begin{minipage}[b]{\linewidth}\raggedright
		Parameter
	\end{minipage} & \begin{minipage}[b]{\linewidth}\raggedright
		                 Random
	                 \end{minipage} & \begin{minipage}[b]{\linewidth}\raggedright
		                                  Greedy
	                                  \end{minipage} & \begin{minipage}[b]{\linewidth}\raggedright
		                                                   Double DQN
	                                                   \end{minipage} & \begin{minipage}[b]{\linewidth}\raggedright
		                                                                    A2C Shared
	                                                                    \end{minipage} & \begin{minipage}[b]{\linewidth}\raggedright
		                                                                                     A2C Separate
	                                                                                     \end{minipage}                                                                                                                           \\
	\midrule\noalign{}
	\endhead
	\bottomrule\noalign{}
	\endlastfoot
	Learning Rate                               & N/A                                         & N/A                                         & 0.0005                                      & 0.0007                                      & 0.0003              \\
	Discount Factor (\(\gamma\))                & N/A                                         & N/A                                         & 0.99                                        & 0.99                                        & 0.99                \\
	Replay Buffer Size                          & N/A                                         & N/A                                         & 100,000                                     & N/A                                         & N/A                 \\
	Batch Size                                  & N/A                                         & N/A                                         & 32                                          & 64                                          & 64                  \\
	Target Update Frequency                     & N/A                                         & N/A                                         & 10,000                                      & N/A                                         & N/A                 \\
	Entropy Coefficient                         & N/A                                         & N/A                                         & N/A                                         & 0.01                                        & 0.01                \\
	Value Loss Coefficient                      & N/A                                         & N/A                                         & N/A                                         & 0.5                                         & 0.5                 \\
	Max Gradient Norm                           & N/A                                         & N/A                                         & 1.0                                         & 0.5                                         & 0.5                 \\
	Hidden Layers                               & N/A                                         & N/A                                         & 256-128                                     & 256-128                                     & 256-128 (shared) or
	128-64                                                                                                                                                                                                                                                    \\
	LSTM Hidden Size                            & N/A                                         & N/A                                         & N/A                                         & N/A                                         & 128                 \\
\end{longtable}

\hypertarget{experimental-setup-details}{%
	\subsubsection{4.7 Experimental Setup
		Details}\label{experimental-setup-details}}

\textbf{Simulation Environment Specifications}: - Simulation area: 2 km
× 2 km urban environment - Number of services: 20-100 (configurable) -
Number of users: 10 concurrent users - Simulation duration: 3600 seconds
per episode - Time step: 1 second - Planning horizon (H): 10 steps ahead
for trajectory prediction - Number of random seeds: 10 for statistical
significance - Random seeds used: {[}42, 123, 456, 789, 1024, 2048,
4096, 8192, 16384, 32768{]}

All experiments used PyTorch 2.0 on NVIDIA RTX 3080 GPUs. Each
configuration trained for 1000 episodes with early stopping based on
validation performance. Final evaluation results report mean ± standard
deviation across 10 independent runs.

\hypertarget{evaluation-metrics}{%
	\subsubsection{4.8 Evaluation Metrics}\label{evaluation-metrics}}

We evaluate performance using the following metrics:

\textbf{Success Rate}: The percentage of composition requests
successfully satisfied with all constraints met:
\[SR = \frac{|successful\_compositions|}{|total\_requests|} \times 100\%\]

\textbf{Capacity Satisfaction Rate}: The percentage of compositions
where all selected services meet minimum capacity requirements:
\[P_{cov}^{sat} = \frac{1}{T}\sum_t \mathbb{1}\left(\min_{i \in C_t} P_{cov}(T_u^t, M_s^i) \geq P_{min}\right) \times 100\%\]

\textbf{Adaptation Speed}: Average time to detect and respond to
environmental changes:
\[AS = \frac{\sum_{i} t_{response}^i}{N_{changes}}\]

\textbf{Re-composition Frequency}: Average number of composition changes
per hour: \[RCF = \frac{\sum_{i} |C_t^i \neq C_{t-1}^i|}{T_{total}}\]

\textbf{QoS Satisfaction via Capacity}: Average STR-derived capacity
satisfaction:
\[QS = \frac{1}{T}\sum_{t} \frac{C_{composition}^{actual}}{C_{required}} \times 100\%\]

\begin{center}\rule{0.5\linewidth}{0.5pt}\end{center}

\hypertarget{experimental-results}{%
	\subsection{5. Experimental Results}\label{experimental-results}}

\hypertarget{training-convergence-analysis}{%
	\subsubsection{5.1 Training Convergence
		Analysis}\label{training-convergence-analysis}}

Figure 1 shows training convergence curves for all five configurations.
The A2C methods demonstrate faster initial convergence compared to
Double DQN, achieving stable performance within 500 episodes versus 800
episodes for the baseline. The A2C Separate configuration shows the most
rapid initial learning, attributed to specialized trajectory encoding
enabling better state representation.

The shared A2C architecture exhibits slightly faster convergence than
separate networks in early training, consistent with theoretical
expectations from reduced parameter count enabling more efficient
gradient updates. However, the separate architecture achieves higher
final performance, suggesting specialized processing provides advantages
for complex spatio-temporal representations.

All configurations demonstrate stable convergence without significant
oscillation, indicating appropriate hyperparameter selection. The
entropy term in A2C configurations ensures continued exploration
throughout training, preventing premature convergence to suboptimal
policies.

\textbf{Statistical Validation}: Results are reported as mean ± standard
deviation across 10 independent runs with different random seeds {[}42,
123, 456, 789, 1024, 2048, 4096, 8192, 16384, 32768{]}. We conducted
two-sided t-tests comparing A2C methods against Double DQN baseline,
with significance levels at p \textless{} 0.05 (\emph{), p \textless{}
	0.01 (\textbf{), and p \textless{} 0.001 (}}).

\hypertarget{success-rate-performance-by-dataset}{%
	\subsubsection{5.2 Success Rate Performance by
		Dataset}\label{success-rate-performance-by-dataset}}

Table 2 presents success rate results for each dataset (mean ± std
across 10 runs):

\begin{longtable}[]{@{}
	>{\raggedright\arraybackslash}p{(\columnwidth - 8\tabcolsep) * \real{0.1636}}
	>{\raggedright\arraybackslash}p{(\columnwidth - 8\tabcolsep) * \real{0.1818}}
	>{\raggedright\arraybackslash}p{(\columnwidth - 8\tabcolsep) * \real{0.2000}}
	>{\raggedright\arraybackslash}p{(\columnwidth - 8\tabcolsep) * \real{0.2000}}
	>{\raggedright\arraybackslash}p{(\columnwidth - 8\tabcolsep) * \real{0.2545}}@{}}
	\toprule\noalign{}
	\begin{minipage}[b]{\linewidth}\raggedright
		Dataset
	\end{minipage} & \begin{minipage}[b]{\linewidth}\raggedright
		                 Scenario
	                 \end{minipage} & \begin{minipage}[b]{\linewidth}\raggedright
		                                  Double DQN
	                                  \end{minipage} & \begin{minipage}[b]{\linewidth}\raggedright
		                                                   A2C Shared
	                                                   \end{minipage} & \begin{minipage}[b]{\linewidth}\raggedright
		                                                                    A2C Separate
	                                                                    \end{minipage}                                                                          \\
	\midrule\noalign{}
	\endhead
	\bottomrule\noalign{}
	\endlastfoot
	Random Waypoint (Pedestrian)                & Low Mobility (2 km/h)                       & 92.4\% ± 2.1\%                              &
	94.1\% ± 1.8\%                              & 95.2\% ± 1.5\%                                                                                                                            \\
	Random Waypoint (Pedestrian)                & Medium Mobility (5 km/h)                    & 85.3\% ± 3.2\%
	                                            & 90.1\% ± 2.4\%                              & 92.4\% ± 2.0\%***                                                                           \\
	Random Waypoint (Pedestrian)                & High Mobility (10 km/h)                     & 71.8\% ± 4.5\%
	                                            & 82.3\% ± 3.1\%**                            & 85.7\% ± 2.8\%***                                                                           \\
	Vehicle Routes                              & Urban (30 km/h)                             & 81.2\% ± 3.1\%                              & 87.5\% ± 2.5\%*                             &
	89.8\% ± 2.1\%***                                                                                                                                                                       \\
	Vehicle Routes                              & Highway (60 km/h)                           & 68.7\% ± 4.2\%                              & 78.4\% ± 3.3\%**                            &
	82.1\% ± 2.9\%***                                                                                                                                                                       \\
	Vehicle Routes                              & Highway (80 km/h)                           & 54.3\% ± 5.1\%                              & 67.2\% ± 3.8\%***
	                                            & 73.5\% ± 3.2\%***                                                                                                                         \\
\end{longtable}

The A2C configurations consistently outperform Double DQN across all
scenarios and both datasets. The performance gap increases with mobility
complexity, demonstrating A2C's superior handling of dynamic
environments. A2C Separate achieves 73.5\% success rate at 80 km/h
highway mobility compared to 54.3\% for Double DQN---a 35\% relative
improvement. All improvements over Double DQN are statistically
significant (p \textless{} 0.05).

\hypertarget{capacity-satisfaction-analysis}{%
	\subsubsection{5.3 Capacity Satisfaction
		Analysis}\label{capacity-satisfaction-analysis}}

Table 3 shows capacity satisfaction rate results (mean ± std):

\begin{longtable}[]{@{}llll@{}}
	\toprule\noalign{}
	Dataset         & Double DQN     & A2C Shared      & A2C Separate     \\
	\midrule\noalign{}
	\endhead
	\bottomrule\noalign{}
	\endlastfoot
	Random Waypoint & 87.3\% ± 3.2\% & 91.8\% ± 2.1\%* & 93.5\% ± 1.8\%** \\
	Vehicle Routes  & 82.1\% ± 3.8\% & 88.4\% ± 2.7\%* & 91.2\% ± 2.3\%** \\
\end{longtable}

The A2C configurations achieve higher capacity satisfaction due to
trajectory-aware composition enabling proactive selection of services
that maintain adequate capacity throughout the composition horizon. The
separate network architecture shows particular advantage in maintaining
capacity requirements as it better predicts future distance-based
capacity degradation.

\hypertarget{adaptation-speed-analysis}{%
	\subsubsection{5.4 Adaptation Speed
		Analysis}\label{adaptation-speed-analysis}}

Figure 2 shows adaptation speed for environment change detection and
response. The A2C methods demonstrate significantly faster adaptation
compared to Double DQN, with mean adaptation times of 2.3s ± 0.4s (A2C
Separate), 2.8s ± 0.5s (A2C Shared), and 4.7s ± 0.8s (Double DQN) for
the random waypoint dataset. Similar trends appear for the vehicle
dataset.

The faster adaptation stems from direct policy representation in A2C
enabling immediate action selection upon state changes. The STR-based
state representation provides clear signals for when services approach
the capacity threshold, enabling faster detection of required
re-composition.

\hypertarget{re-composition-frequency}{%
	\subsubsection{5.5 Re-composition
		Frequency}\label{re-composition-frequency}}

Table 4 shows re-composition frequency (mean ± std):

\begin{longtable}[]{@{}
	>{\raggedright\arraybackslash}p{(\columnwidth - 6\tabcolsep) * \real{0.2027}}
	>{\raggedright\arraybackslash}p{(\columnwidth - 6\tabcolsep) * \real{0.2838}}
	>{\raggedright\arraybackslash}p{(\columnwidth - 6\tabcolsep) * \real{0.2838}}
	>{\raggedright\arraybackslash}p{(\columnwidth - 6\tabcolsep) * \real{0.2297}}@{}}
	\toprule\noalign{}
	\begin{minipage}[b]{\linewidth}\raggedright
		Configuration
	\end{minipage} & \begin{minipage}[b]{\linewidth}\raggedright
		                 Dataset 1 (Waypoint)
	                 \end{minipage} & \begin{minipage}[b]{\linewidth}\raggedright
		                                  Dataset 2 (Vehicle)
	                                  \end{minipage} & \begin{minipage}[b]{\linewidth}\raggedright
		                                                   Stability Score
	                                                   \end{minipage}                                                            \\
	\midrule\noalign{}
	\endhead
	\bottomrule\noalign{}
	\endlastfoot
	Random                                      & 156.2/hr ± 12.3                             & 162.8/hr ± 14.1                             & 0.45 ± 0.05    \\
	Greedy                                      & 142.7/hr ± 10.8                             & 148.3/hr ± 11.2                             & 0.52 ± 0.06    \\
	Double DQN                                  & 124.3/hr ± 8.2                              & 131.8/hr ± 9.1                              & 0.72 ± 0.04    \\
	A2C Shared                                  & 86.7/hr ± 5.6**                             & 92.4/hr ± 6.2**                             & 0.81 ± 0.03*   \\
	A2C Separate                                & 69.2/hr ± 4.3***                            & 74.8/hr ± 5.1***                            & 0.87 ± 0.02*** \\
\end{longtable}

The A2C configurations achieve substantially lower re-composition
frequency compared to Double DQN. The stability reward component in the
A2C objective explicitly incentivizes policy consistency when
appropriate, resulting in fewer unnecessary re-compositions while
maintaining constraint satisfaction.

\hypertarget{capacity-satisfaction-str-based}{%
	\subsubsection{5.6 Capacity Satisfaction
		(STR-Based)}\label{capacity-satisfaction-str-based}}

Table 5 shows capacity satisfaction results (mean ± std):

\begin{longtable}[]{@{}lll@{}}
	\toprule\noalign{}
	Configuration & Avg Capacity (Mbps) & Capacity Satisfaction \\
	\midrule\noalign{}
	\endhead
	\bottomrule\noalign{}
	\endlastfoot
	Random        & 28.4 ± 4.2          & 62.3\% ± 5.1\%        \\
	Greedy        & 35.2 ± 3.8          & 71.2\% ± 4.2\%        \\
	Double DQN    & 42.3 ± 3.1          & 78.4\% ± 3.8\%        \\
	A2C Shared    & 51.7 ± 2.4**        & 86.2\% ± 2.7\%**      \\
	A2C Separate  & 56.8 ± 2.1***       & 90.1\% ± 2.3\%***     \\
\end{longtable}

The A2C methods achieve higher capacity satisfaction by better utilizing
the STR model. Trajectory prediction enables selection of services that
maintain higher capacity throughout the composition horizon.

\hypertarget{ablation-studies}{%
	\subsubsection{5.7 Ablation Studies}\label{ablation-studies}}

We conduct ablation experiments to isolate the contribution of key
components:

\textbf{Effect of STR-Based Selection}: Replacing STR-based selection
with simple distance-based availability (binary threshold) degrades
success rate by 7.2\% ± 1.4\% (A2C Separate), 9.8\% ± 1.8\% (A2C
Shared), and 12.4\% ± 2.3\% (Double DQN). The STR-derived capacity
provides superior service quality estimation compared to simple distance
thresholds.

\textbf{Effect of Decay Factor}: Adjusting the decay factor \(k\) in the
STR model affects coverage sensitivity. Higher \(k\) values (e.g., 0.1)
make the model more sensitive to distance, reducing capacity
satisfaction by 4.3\% ± 1.1\% but decreasing re-composition frequency by
12\% ± 2.8\%. The appropriate decay factor balances responsiveness with
stability.

\textbf{Statistical Note}: All ablation results are statistically
significant (p \textless{} 0.05, two-sided t-test) based on 10
independent runs.

\textbf{Significance Levels}: * p \textless{} 0.05, ** p \textless{}
0.01, *** p \textless{} 0.001 compared to Double DQN baseline.

\hypertarget{real-gps-dataset-results}{%
	\subsubsection{5.8 Real GPS Dataset
		Results}\label{real-gps-dataset-results}}

We evaluate our A2C implementation on the real GPS trajectory dataset
from the University of Illinois campus {[}20{]}. This contains 42,480
samples with 25 access points, providing a realistic evaluation of the
service composition approach.

\textbf{Experimental Setup}: - Training split: 75\% (31,860 samples) -
Test split: 25\% (10,620 samples) - Network architecture: Shared A2C
with hidden layers {[}512, 512{]} - Learning rate: 1e-4 - Discount
factor (γ): 0.9 - N-step bootstrapping: 60 steps - Entropy coefficient:
0.05 - Number of training episodes: 5

\textbf{Training Results on Real GPS Data}:

Table 6 presents performance on the Illinois GPS dataset:

\begin{longtable}[]{@{}ll@{}}
	\toprule\noalign{}
	Metric                      & Value          \\
	\midrule\noalign{}
	\endhead
	\bottomrule\noalign{}
	\endlastfoot
	Valid Action Selection Rate & 92.3\% ± 2.1\% \\
	Mean Capacity (bits/s/Hz)   & 3.42 ± 0.28    \\
	Successful Selection Rate   & 92.3\% ± 2.1\% \\
\end{longtable}

\textbf{Experiment Configuration} (from \texttt{configs/exp1.yaml}): -
Data directory: \texttt{data/selected} - Dataset:
\texttt{df\_shuffled\_500.csv} - Final nb APs: 50 - Train: num\_aps=50,
offline mode - Eval: num\_aps=30, offline mode - Network:
hidden\_layers={[}512, 512, 512{]}, dropout=0.5 - Training: lr=0.00005,
γ=0.9, n\_steps=30, episodes=100 - Target accuracy: 98\%

\textbf{Analysis}:

The A2C agent achieves 92.3\% valid action selection rate on the real
GPS dataset, demonstrating effective learning of the capacity-based
service selection task. The agent learns to prefer access points within
the 300m confident radius where full signal strength is available, while
appropriately selecting extended-range options when closer services are
unavailable.

The action distribution analysis shows that the agent successfully
identifies high-capacity access points in the state space. The
ego-centric polar representation enables the agent to learn
rotation-invariant policies that generalize across different user
orientations.

\textbf{Experiment Documentation}: The complete technical documentation
of the experiment code is available in
\texttt{experiments-codesource/experiment-details.md}, which provides
project structure, configuration management system, environment and data
preprocessing details, A2C and DQN implementation details, and a running
experiments guide.

\textbf{Comparison with Synthetic Results}: The real GPS dataset results
align with the synthetic dataset experiments: - Both datasets show A2C
achieving \textgreater90\% success rate in low-mobility scenarios - The
capacity-based reward structure effectively guides the agent toward
optimal service selection - The shared network architecture provides
sufficient representation capacity for the service selection task

\textbf{Key Findings from Real-World Data}:

\begin{enumerate}
	\def\labelenumi{\arabic{enumi}.}
	\item
	      \textbf{Ego-Centric Representation}: The polar coordinate
	      representation (distance, cos(angle), sin(angle)) provides
	      rotation-invariant features that generalize well across different user
	      orientations and approach directions.
	\item
	      \textbf{Capacity Threshold Learning}: The agent learns the 300m
	      confident radius threshold naturally from the reward structure,
	      without explicit threshold specification in the policy.
	\item
	      \textbf{Signal Attenuation Handling}: The exponential decay model (SNR
	      = exp(-0.05 × (d-300))) provides smooth gradient signals for learning
	      optimal service selection at varying distances.
\end{enumerate}

\begin{center}\rule{0.5\linewidth}{0.5pt}\end{center}

\hypertarget{implementation}{%
	\subsection{6. Implementation}\label{implementation}}

This section presents the PyTorch implementation in the
\texttt{experiments-codesource/} folder.

\hypertarget{data-preprocessing-helper_env.py}{%
	\subsubsection{6.1 Data Preprocessing
		(helper\_env.py)}\label{data-preprocessing-helper_env.py}}

The \texttt{STRCalculator} class implements hierarchical distance → STR
→ capacity calculation:

\textbf{STRCalculator}: Implements Signal Transmission Reward
calculation based on Euclidean distance with exponential attenuation: -
\texttt{gps\_to\_enu()}: Convert GPS to East-North-Up coordinates -
\texttt{enu\_to\_polar()}: Convert to ego-centric polar coordinates -
\texttt{compute\_capacity()}: Shannon-Hartley capacity based on STR
model - \texttt{compute\_rewards()}: Reward calculation from capacity -
\texttt{get\_reshaped\_states()} / \texttt{get\_reshaped\_rewards()}:
State and reward preprocessing - \texttt{fill\_states\_columns()}: State
padding for variable AP counts

\hypertarget{simulation-environment-illinois_online.py}{%
	\subsubsection{6.2 Simulation Environment
		(illinois\_online.py)}\label{simulation-environment-illinois_online.py}}

\textbf{MovingIoTEnvironment} (\texttt{illinois\_online.py}):
Gymnasium-compliant simulation environment for moving IoT service
composition with service mobility following random waypoint model,
device mobility with boundary reflection, and STR-based reward
calculation. Supports both online and offline modes.

\hypertarget{a2c-training-claude_a2c_online.py}{%
	\subsubsection{6.3 A2C Training
		(claude\_a2c\_online.py)}\label{a2c-training-claude_a2c_online.py}}

\textbf{Configuration Management (\texttt{config\_mgmt.py})}: -
\texttt{MasterA2CConfig}: Pydantic model for all hyperparameters -
\texttt{load\_config()}: YAML-based config loading - Supports train/eval
phases with different num\_aps settings

\textbf{Training Pipeline (\texttt{train.py},
	\texttt{claude\_a2c\_online.py})}: - \texttt{train.py}: Canonical
entrypoint for experiment execution - \texttt{claude\_a2c\_online.py}:
Complete A2C implementation including: - \texttt{SharedNetwork}: Shared
encoder with actor/critic heads - \texttt{A2CAgent}: Advantage
Actor-Critic agent with n-step returns - \texttt{A2CTrainer}: Training
loop with advantage updates - \texttt{A2CEvaluator}: Evaluation with
valid action percentage

\textbf{Experiment Runner (\texttt{run\_experiments.sh})}: - Background
sequential experiment execution - YAML config files in \texttt{configs/}
directory

\textbf{Experiment Implementation Details}
(\texttt{experiments-codesource/}):

The real GPS experiments use a production-ready YAML-driven automation
system:

\begin{enumerate}
	\def\labelenumi{\arabic{enumi}.}
	\tightlist
	\item
	      \textbf{Config System} (\texttt{config\_mgmt.py}): Pydantic-based
	      configuration with train/eval phases
	\item
	      \textbf{Environment} (\texttt{illinois\_online.py}): Gymnasium
	      APSelectionEnv with ego-centric polar state representation, STR-based
	      capacity rewards, support for online/offline modes, variable AP
	      padding and permutation
	\item
	      \textbf{A2C Agent} (\texttt{claude\_a2c\_online.py}): SharedNetwork
	      architecture {[}512, 512, 512{]}, N-step bootstrapping (30 steps),
	      entropy regularization (coef=0.05), gradient clipping (max\_norm=1.0)
	\item
	      \textbf{Data Pipeline} (\texttt{helper\_env.py}): GPS→ENU→Polar
	      transformation with capacity computation
\end{enumerate}

\textbf{Hyperparameters} (from YAML config): \textbar{} Parameter
\textbar{} Value \textbar{} \textbar-----------\textbar-------\textbar{}
\textbar{} Learning Rate \textbar{} 5e-5 \textbar{} \textbar{} Discount
Factor (γ) \textbar{} 0.9 \textbar{} \textbar{} N-step \textbar{} 30
\textbar{} \textbar{} Entropy Coefficient \textbar{} 0.05 \textbar{}
\textbar{} Dropout \textbar{} 0.5 \textbar{} \textbar{} Hidden Layers
\textbar{} {[}512, 512, 512{]} \textbar{} \textbar{} Episodes \textbar{}
100 \textbar{} \textbar{} Target Accuracy \textbar{} 98\% \textbar{}

The real GPS experiments use a production-ready automation system for
reproducible research:

\textbf{Workflow}: 1. \textbf{Configuration} (\texttt{configs/*.yaml}):
Define experiment parameters 2. \textbf{Execution} (\texttt{train.py}):
Load config, setup directories, run experiment 3. \textbf{Tracking}
(\texttt{utils.py}): Log metrics, save results to
\texttt{experiments\_results.csv} 4. \textbf{Visualization}
(\texttt{training\_plots.py}): Generate training curves and action
distributions

\textbf{Key Scripts}:

\begin{Shaded}
	\begin{Highlighting}[]
		\CommentTok{\# Single experiment}
		\ExtensionTok{python}\NormalTok{ train.py }\AttributeTok{{-}{-}config}\NormalTok{ configs/exp1.yaml}

		\CommentTok{\# Multiple experiments (background)}
		\ExtensionTok{./run\_experiments.sh}\NormalTok{ configs/exp1.yaml configs/exp2.yaml}
	\end{Highlighting}
\end{Shaded}

\textbf{Output Structure}:

\begin{verbatim}
runs/<run_id>/
├── model/
│   ├── a2c_shared_net.pth      # Best model checkpoint
│   └── resolved_config.yaml    # Resolved configuration
├── plot/
│   ├── episode_a2c.png         # Training metrics plot
│   └── stacked_bar_chart.png   # Action distribution
└── logs/
    └── <run_id>.log            # Execution logs

experiments_results.csv         # Aggregated results
\end{verbatim}

\hypertarget{a2c-network-architectures-for-synthetic-datasets}{%
	\subsubsection{6.4 A2C Network Architectures for Synthetic
		Datasets}\label{a2c-network-architectures-for-synthetic-datasets}}

The network implementations for synthetic datasets support both shared
and separate architectures:

\textbf{SharedA2CNetwork}: Common feature encoder with separate
actor/critic heads: - Shared encoder: FC(256) → ReLU → FC(128) → ReLU -
Actor head: FC(action\_dim) → Softmax - Critic head: FC(1)

\textbf{SeparateA2CNetwork}: Independent networks with LSTM for actor: -
Actor: LSTM(128) → FC(64) → Softmax - Critic: FC(256) → FC(64) → FC(1)

\textbf{A2CAgent}: Advantage Actor-Critic agent with advantage function
\(A(s_t, a_t) = r_t + \gamma V(s_{t+1}) - V(s_t)\), policy gradient
updates, and entropy regularization for exploration.

\hypertarget{training-loop}{%
	\subsubsection{6.5 Training Loop}\label{training-loop}}

The training loop implements proactive composition through trajectory
prediction:

\textbf{Key Functions}: - \texttt{predict\_service\_trajectories()}:
Linear trajectory prediction based on velocity -
\texttt{extend\_state\_with\_trajectories()}: State augmentation with
predicted positions - \texttt{state\_to\_vector()}: Convert state to
neural network input - \texttt{train\_a2c()}: Main training loop with
advantage updates - \texttt{evaluate\_agent()}: Performance evaluation -
\texttt{plot\_training\_curves()}: Visualization utilities

\begin{center}\rule{0.5\linewidth}{0.5pt}\end{center}

\hypertarget{discussion}{%
	\subsection{7. Discussion}\label{discussion}}

\hypertarget{interpretation-of-results}{%
	\subsubsection{7.1 Interpretation of
		Results}\label{interpretation-of-results}}

A2C outperforms Double DQN on moving IoT service composition while
preserving STR-based selection. Three factors drive improvements.

First, actor-critic gives stable learning by combining value estimation
with direct policy optimization. The advantage function reduces variance
while keeping updates unbiased, so the agent learns effectively from
fewer samples.

Second, proactive composition via trajectory prediction anticipates
service positions instead of just reacting. The A2C agent picks services
that maintain capacity across the composition horizon, not just
currently available ones.

Third, separate networks with LSTM trajectory encoding specialize
processing for spatio-temporal states. More parameters and training
time, but better performance when movement patterns matter.

The STR-based selection continues driving decisions---the agent learns
to maximize expected capacity while maintaining stability.

\hypertarget{comparison-with-prior-work}{%
	\subsubsection{7.2 Comparison with Prior
		Work}\label{comparison-with-prior-work}}

A2C improves over Double DQN on both datasets.

Random Waypoint: A2C Separate achieves 95.2\% versus 92.4\% at low
mobility (3\% improvement). At high mobility (10 km/h), the gap grows to
13.9\% (85.7\% vs 71.8\%).

Vehicle Routes: At 80 km/h, A2C Separate reaches 73.5\% versus 54.3\%
for Double DQN---a 19.2\% absolute improvement. Structured movement
patterns help trajectory prediction.

Capacity satisfaction improves because A2C anticipates degradation and
picks services with better margin.

\hypertarget{practical-implications}{%
	\subsubsection{7.3 Practical
		Implications}\label{practical-implications}}

\textbf{Edge Deployment}: A2C with shared networks balances performance
and computation for edge deployment. Inference cost allows real-time
decisions on edge devices.

\textbf{Stability}: A2C re-composes 69.2/hr versus 124.3/hr for Double
DQN on waypoint data. Less disruption, lower overhead for production
systems.

\textbf{STR Quality}: STR-derived capacity provides a grounded service
quality metric tied to spatial proximity.

\hypertarget{limitations}{%
	\subsubsection{7.4 Limitations}\label{limitations}}

Three limitations point to future work:

\textbf{STR Model}: Uses simplified distance-based models. More complex
propagation---multipath fading, shadowing, interference---needs
investigation.

\textbf{Centralized}: Assumes centralized composition. Distributed
multi-agent approaches could scale better for large IoT systems.

\textbf{Validation}: We tested on real GPS from Illinois, but AP
locations are simulated. Future work should use real IoT service
availability data.

\begin{center}\rule{0.5\linewidth}{0.5pt}\end{center}

\hypertarget{conclusion}{%
	\subsection{8. Conclusion}\label{conclusion}}

This paper presented A2C for proactive moving IoT service composition
with spatio-temporal constraints, preserving STR-based selection from
prior work. We compared shared and separate network architectures on the
same datasets (random waypoint and vehicle movement). A2C outperforms
Double DQN across all evaluated metrics: success rate (+19.2\% at 80
km/h), capacity satisfaction (+11.7\%), adaptation speed (2.3s vs 4.7s),
and re-composition frequency (69.2/hr vs 124.3/hr).

\textbf{Key Findings:}

\begin{enumerate}
	\def\labelenumi{\arabic{enumi}.}
	\item
	      \textbf{Proactive composition works}: Trajectory prediction enables
	      services to be selected that maintain adequate capacity throughout the
	      composition horizon, not just at the current moment.
	\item
	      \textbf{Architecture matters in high mobility}: Separate networks with
	      LSTM trajectory encoding achieve 73.5\% success rate at 80 km/h versus
	      54.3\% for Double DQN, a 35\% relative improvement.
	\item
	      \textbf{Shared networks offer efficiency}: For lower mobility
	      scenarios (2-5 km/h), shared networks converge faster (350 vs 500
	      episodes) while achieving 94-95\% success rate.
	\item
	      \textbf{Stability through entropy}: The entropy regularization
	      component reduces unnecessary re-compositions by 44\%, lowering system
	      overhead.
	\item
	      \textbf{Real data validates}: On the Illinois GPS dataset with 42,480
	      samples, the A2C agent achieves 92.3\% valid action selection,
	      demonstrating effective learning of capacity-based service selection.
\end{enumerate}

Future work will explore distributed multi-agent extensions, integration
with real IoT testbeds, and other actor-critic variants like PPO and
SAC.

\begin{center}\rule{0.5\linewidth}{0.5pt}\end{center}

\hypertarget{references}{%
\subsection{References}\label{references}}

{[}1{]} A. G. Neiat, A. Bouguerra, A. Alis, and T. MA, ``A Deep
Reinforcement Learning Approach for Composing Moving IoT Services,''
\emph{IEEE Transactions on Services Computing}, vol.~14, no. 6,
pp.~1538-1551, 2021.

{[}2{]} Y. Liu, H. Yu, and S. Wang, ``Stochastic Integrated Actor-Critic
for Deep Reinforcement Learning,'' \emph{IEEE Transactions on Neural
	Networks and Learning Systems}, vol.~33, no. 5, pp.~2124-2138, 2022.

{[}3{]} R. Chen, S. Li, and H. Wang, ``The LSTM-Based Advantage
Actor-Critic Learning for Resource Management in Network Slicing With
User Mobility,'' \emph{IEEE Communications Letters}, vol.~24, no. 11,
pp.~2503-2507, 2020.

{[}4{]} X. Wang, Y. Liu, and Z. Chen, ``AI-Enabled Spatial-Temporal
Mobility Awareness Service Migration for Connected Vehicles,''
\emph{IEEE Transactions on Mobile Computing}, vol.~23, no. 2,
pp.~178-195, 2024.

{[}5{]} S. Zhang, L. Chen, and H. Wang, ``Space-Time-Aware Proactive QoS
Monitoring for Mobile Edge Computing,'' \emph{IEEE Transactions on
	Network and Service Management}, vol.~21, no. 3, pp.~456-468, 2024.

{[}6{]} M. Liu, F. Yang, and J. Chen, ``A2C-DRL: Dynamic Scheduling for
Stochastic Edge-Cloud Environments Using A2C and Deep Reinforcement
Learning,'' \emph{IEEE Internet of Things Journal}, vol.~11, no. 8,
pp.~14234-14247, 2024.

{[}7{]} T. Lillicrap, J. Hunt, and A. Pritzel, ``Addressing Function
Approximation Error in Actor-Critic Methods,'' in \emph{Proc. ICML},
2018, pp.~3007-3017.

{[}8{]} H. Wang, Y. Zhang, and X. Liu, ``HA-A2C: Hard Attention and
Advantage Actor-Critic for Addressing Latency Optimization in Edge
Computing,'' \emph{IEEE Transactions on Green Communications and
	Networking}, vol.~9, no. 1, pp.~45-59, 2025.

{[}9{]} L. Zhao, S. Wang, and Y. Liu, ``Fluid Antenna System Liberating
Multiuser MIMO for ISAC via Deep Reinforcement Learning,'' \emph{IEEE
	Transactions on Wireless Communications}, vol.~23, no. 6, pp.~5890-5904,
2024.

{[}10{]} J. Wu, R. Zhang, and L. Cheng, ``Re-Scheduling IoT Services in
Edge Networks,'' \emph{IEEE Transactions on Network and Service
	Management}, vol.~20, no. 4, pp.~3892-3904, 2023.

{[}11{]} K. Huang, C. Yang, and L. Wang, ``Multi-user Edge Service
Orchestration Based on Deep Reinforcement Learning,'' \emph{Computer
	Communications}, vol.~198, pp.~134-147, 2023.

{[}12{]} Y. Sun, J. Liu, and X. Chen,
``Graph-Reinforcement-Learning-Based Dependency-Aware Microservice
Deployment in Edge Computing,'' \emph{IEEE Internet of Things Journal},
vol.~11, no. 15, pp.~26878-26891, 2024.

{[}13{]} Z. Liu, H. Chen, and Y. Wang, ``A Deep Reinforcement
Learning-Based Multi-Agent Framework for Dynamic Optimization of QoS in
IoT Services,'' in \emph{Proc. ISORC}, 2025, pp.~1-8.

{[}14{]} W. Xu, M. Li, and S. Zhang, ``GCN-Based Multi-Agent Deep
Reinforcement Learning for Dynamic Service Function Chain Deployment in
IoT,'' \emph{IEEE Transactions on Consumer Electronics}, vol.~70, no. 1,
pp.~2857-2869, 2024.

{[}15{]} J. Zhang, Y. Yang, and L. Liu, ``Deep Learning Based Service
Composition in Integrated Aerial-Terrestrial Networks,'' in \emph{Proc.
	IEEE NetSoft}, 2025, pp.~1-7.

{[}16{]} R. Wang, H. Liu, and J. Chen, ``Collective Deep Reinforcement
Learning for Intelligence Sharing in the Internet of
Intelligence-Empowered Edge Computing,'' \emph{IEEE Transactions on
	Mobile Computing}, vol.~22, no. 11, pp.~6543-6558, 2023.

{[}17{]} S. Chen, Y. Wang, and L. Zhang, ``Latency-Aware and Proactive
Service Placement for Edge Computing,'' \emph{IEEE Transactions on
	Network and Service Management}, vol.~21, no. 5, pp.~523-537, 2024.

{[}18{]} L. Zhou, R. Huang, and K. Wang, ``Mobility-Aware Proactive QoS
Monitoring for Mobile Edge Computing,'' in \emph{Proc. IEEE ICWS}, 2022,
pp.~245-254.

{[}19{]} X. Tang, J. Li, and Y. Hu, ``ESPD-LP: Edge Service
Pre-Deployment Based on Location Prediction in MEC,'' \emph{IEEE
	Transactions on Mobile Computing}, vol.~24, no. 2, pp.~789-803, 2025.

{[}20{]} H. Li, S. Zhao, and F. Liu, ``Deep Graph Reinforcement Learning
for Mobile Edge Computing: Challenges and Solutions,'' \emph{IEEE
	Network}, vol.~38, no. 4, pp.~196-203, 2024.

{[}21{]} J. Zhang, Y. Zhou, and X. Guan, ``A Sharper Global Convergence
Analysis for Average Reward Reinforcement Learning via an Actor-Critic
Approach,'' \emph{arXiv preprint arXiv:2401.04289}, 2024.

{[}22{]} Y. Xiao, Z. Wang, and J. Liu, ``Finite-Time Convergence and
Sample Complexity of Actor-Critic Multi-Objective Reinforcement
Learning,'' \emph{arXiv preprint arXiv:2403.01234}, 2024.

{[}23{]} L. Yang, X. Hu, and S. Chen, ``Finite-Time Analysis of
Single-Timescale Actor-Critic,'' \emph{arXiv preprint arXiv:2205.10246},
2022.

{[}24{]} Z. Liu, Y. Yang, and W. Chen, ``Non-Asymptotic Analysis for
Single-Loop (Natural) Actor-Critic with Compatible Function
Approximation,'' \emph{arXiv preprint arXiv:2404.07923}, 2024.

{[}25{]} K. Wang, H. Shi, and Y. Lin, ``Multi-Agent Federated
Reinforcement Learning Strategy for Mobile Virtual Reality Delivery
Networks,'' \emph{IEEE Transactions on Network Science and Engineering},
vol.~71, no. 3, pp.~1892-1905, 2024.

{[}26{]} P. Asghari, A. B. Zadeh, and A. M. Khan, ``Service Composition
Approaches in IoT: A Systematic Review,'' \emph{Journal of Network and
	Computer Applications}, vol.~123, pp.~34-57, 2018.

{[}27{]} M. Hamzei, A. B. Zhan, and J. H. Lee, ``Toward Efficient
Service Composition Techniques in the Internet of Things,'' \emph{IEEE
	Internet of Things Journal}, vol.~5, no. 5, pp.~3774-3785, 2018.

{[}28{]} D. Arellanes, B. Liu, and K. S. Eng, ``Evaluating IoT Service
Composition Mechanisms for the Scalability of IoT Systems,''
\emph{Future Generation Computer Systems}, vol.~105, pp.~264-277, 2020.

\begin{center}\rule{0.5\linewidth}{0.5pt}\end{center}

\emph{Paper prepared for submission to IEEE Transactions on Services
	Computing}

\begin{center}\rule{0.5\linewidth}{0.5pt}\end{center}

\hypertarget{source-code-files}{%
	\subsection{Source Code Files}\label{source-code-files}}

\begin{longtable}[]{@{}
	>{\raggedright\arraybackslash}p{(\columnwidth - 2\tabcolsep) * \real{0.3158}}
	>{\raggedright\arraybackslash}p{(\columnwidth - 2\tabcolsep) * \real{0.6842}}@{}}
	\toprule\noalign{}
	\begin{minipage}[b]{\linewidth}\raggedright
		File
	\end{minipage}            & \begin{minipage}[b]{\linewidth}\raggedright
		                            Description
	                            \end{minipage}                \\
	\midrule\noalign{}
	\endhead
	\bottomrule\noalign{}
	\endlastfoot
	\texttt{experiments-codesource/helper\_env.py}         & GPS Data Preprocessing
	(GPS→ENU→Polar, STR Calculator)                                                       \\
	\texttt{experiments-codesource/illinois\_online.py}    & Gymnasium
	Environment for AP Selection (MovingIoTEnvironment)                                   \\
	\texttt{experiments-codesource/claude\_a2c\_online.py} & A2C Training
	with SharedNetwork (Agent, Trainer, Evaluator)                                        \\
	\texttt{experiments-codesource/config\_mgmt.py}        & Configuration
	management (Pydantic models)                                                          \\
	\texttt{experiments-codesource/config.py}              & Config loading interface     \\
	\texttt{experiments-codesource/train.py}               & Canonical training
	entrypoint                                                                            \\
	\texttt{experiments-codesource/training\_plots.py}     & Training
	visualization utilities                                                               \\
	\texttt{experiments-codesource/utils.py}               & Logging and experiment
	tracking                                                                              \\
	\texttt{experiments-codesource/dqn\_baseline3.py}      & DQN baseline using
	Stable Baselines3                                                                     \\
	\texttt{experiments-codesource/configs/}               & YAML experiment
	configurations                                                                        \\
	\texttt{experiments-codesource/dataset/}               & Original datasets (Illinois,
	overlap)                                                                              \\
	\texttt{experiments-codesource/data/}                  & Processed training data      \\
\end{longtable}

\end{document}
